\documentclass[10pt,a4paper]{article}
\usepackage[top=12mm,bottom=12mm,left=14mm,right=14mm,headsep=2mm,footskip=7mm]{geometry}
\usepackage{fontspec}
\usepackage{xeCJK}
\setmainfont{Noto Sans}
\setsansfont{Noto Sans}
\setCJKmainfont{Noto Sans CJK SC}
\setCJKsansfont{Noto Sans CJK SC}
\usepackage{xcolor}
\usepackage{graphicx}
\usepackage{tikz}
\usetikzlibrary{calc,arrows.meta,positioning}
\usepackage{fancyhdr}
\usepackage{hyperref}
\usepackage{ragged2e}
\usepackage{microtype}
\definecolor{P2Paper}{HTML}{FCFAF7}
\definecolor{P2Ink}{HTML}{2D3238}
\definecolor{P2Muted}{HTML}{70757A}
\definecolor{P2BlueLight}{HTML}{D3EAF5}
\definecolor{P2Apricot}{HTML}{F5D9B8}
\definecolor{P2Rose}{HTML}{E8C7D3}
\definecolor{P2Violet}{HTML}{D7D3EA}
\definecolor{P2Blue}{HTML}{5B9FBE}
\definecolor{P2ApricotAccent}{HTML}{D49757}
\definecolor{P2RoseAccent}{HTML}{B97589}
\definecolor{P2VioletAccent}{HTML}{8177A8}
\definecolor{P2Rule}{HTML}{DDE3E7}
\definecolor{P2Panel}{HTML}{FFFDFC}
\pagecolor{P2Paper}
\color{P2Ink}
\setlength{\parindent}{0pt}
\setlength{\parskip}{3.5pt}
\linespread{1.04}
\hypersetup{hidelinks}
\pagestyle{fancy}
\fancyhf{}
\renewcommand{\headrulewidth}{0pt}
\fancyfoot[L]{\fontsize{6.6}{8}\selectfont\color{P2Muted}PROCESS REPORT · WORKFLOW CONSTRUCTION}
\fancyfoot[R]{\fontsize{6.6}{8}\selectfont\color{P2Muted}\thepage}
\newcommand{\eyebrow}[1]{{\fontsize{7.8}{9.4}\selectfont\bfseries\color{P2RoseAccent}#1}\par\vspace{1pt}}
\newcommand{\pagetitle}[1]{{\fontsize{18.5}{21.5}\selectfont\bfseries\color{P2Ink}#1}\par\vspace{3pt}}
\newcommand{\deck}[1]{{\fontsize{9.0}{13.0}\selectfont\color{P2Ink}\RaggedRight #1}\par\vspace{4pt}}
\newcommand{\tracehead}[2]{{\fontsize{7.1}{8.4}\selectfont\bfseries\color{#1}#2}\par}
\newcommand{\tracetitle}[1]{{\fontsize{10.2}{12.2}\selectfont\bfseries\color{P2Ink}#1}\par\vspace{1.5pt}}
\newcommand{\tracebody}[1]{{\fontsize{8.0}{11.2}\selectfont\color{P2Ink}\RaggedRight #1}\par}
\newcommand{\provenance}[1]{\vspace{2pt}{\color{P2Rule}\hrule height .35pt}\vspace{2.5pt}{\fontsize{6.9}{9.2}\selectfont\color{P2Muted}\RaggedRight\textbf{Original Evidence.} #1}\par}
\tikzset{
 userA/.style={rounded corners=1mm,draw=P2ApricotAccent,line width=.55pt,fill=P2Apricot,fill opacity=.34},
 userR/.style={rounded corners=1mm,draw=P2RoseAccent,line width=.55pt,fill=P2Rose,fill opacity=.34},
 gptB/.style={rounded corners=1mm,draw=P2Blue,line width=.55pt,fill=P2BlueLight,fill opacity=.27},
 gptV/.style={rounded corners=1mm,draw=P2VioletAccent,line width=.55pt,fill=P2Violet,fill opacity=.27},
 arrA/.style={-{Stealth[length=2.0mm,width=1.45mm]},draw=P2RoseAccent,line width=.72pt},
 arrB/.style={-{Stealth[length=2.0mm,width=1.45mm]},draw=P2Blue,line width=.72pt},
 arrV/.style={-{Stealth[length=2.0mm,width=1.45mm]},draw=P2VioletAccent,line width=.72pt},
 tagA/.style={circle,minimum size=4.8mm,inner sep=0pt,draw=P2ApricotAccent,line width=.6pt,fill=P2Paper,font=\fontsize{6.7}{7}\selectfont\bfseries},
 tagR/.style={circle,minimum size=4.8mm,inner sep=0pt,draw=P2RoseAccent,line width=.6pt,fill=P2Paper,font=\fontsize{6.7}{7}\selectfont\bfseries},
 tagB/.style={circle,minimum size=4.8mm,inner sep=0pt,draw=P2Blue,line width=.6pt,fill=P2Paper,font=\fontsize{6.7}{7}\selectfont\bfseries}
}

\begin{document}

\eyebrow{WORKFLOW CONSTRUCTION · 29}
\pagetitle{Local Diagram解决局部Decision，Global Map只做研究故事总览}
\deck{流程图类型越来越多以后，我没有继续把所有内容硬塞进一张“大演变图”。方法选择、参数决策、AI批判各自用最适合的局部图；等这些Decision基本稳定后，再用Global Map把主要研究故事串起来。}
\begin{tikzpicture}[x=1mm,y=1mm]
  \node[anchor=north west,inner sep=0] (g) at (0,0) {\includegraphics[width=52mm]{assets/gpt_local_global.png}};
  \node[anchor=north west,text width=118mm,align=left] at (59,-2) {
    \tracehead{P2Blue}{LOCAL DIAGRAMS}\tracetitle{每个Decision选择最合适的表达语言}
    \tracebody{任务理解、方法形成、参数裁决和AI批判不需要画成同一种图。局部图先把当前问题说明白，再考虑整份报告是否统一。}
    \vspace{8mm}
    \tracehead{P2VioletAccent}{GLOBAL OVERVIEW}\tracetitle{全局图是导航，不是替代}
    \tracebody{Global Decision Evolution Map只保留阶段、核心Decision和方向，把局部图里的关键节点串起来。内容复杂时宁可“总览 + 详细页”，也不把所有细节塞在一张图里。}
    \vspace{8mm}
    \tracehead{P2RoseAccent}{NATIVE AUTHORING}\tracetitle{正式流程图保留可编辑原生源}
    \tracebody{正式流程图继续用draw.io、Figma、Graphviz或TikZ等可编辑工具制作，不让生成式图片替代真实研究结构。}
  };
\end{tikzpicture}
\provenance{该页来自对“多张流程图如何组织、全局演变图如何定位、流程图用什么工具制作”的真实讨论。}
\newpage

\eyebrow{WORKFLOW CONSTRUCTION · 30}
\pagetitle{Task 1开始前，我把Process Report的基本做法先定下来}
\deck{在真正进入Task 1之前，我已经把过程文档怎么做基本定清楚了：什么算一个Decision、Evidence怎么保存、最终确认怎么记录、流程图怎么选、LaTeX怎么排版。接下来新增的内容不再靠“设计”，而要来自真实Task 1研究过程。}
\begin{tikzpicture}[x=1mm,y=1mm]
  \node[anchor=north west,inner sep=0] (g) at (0,0) {\includegraphics[width=50mm]{assets/gpt_final_structure.png}};
  \node[anchor=north west,inner sep=0] (u) at (55,-76) {\includegraphics[width=38mm]{assets/user_final_confirm.png}};
  \coordinate (ga) at ($(g.south east)+(-6mm,14mm)$);
  \coordinate (ua) at ($(u.north west)+(2mm,-4mm)$);
  \draw[arrV] (ga) .. controls (72,-70) and (80,-75) .. (ua);
  \node[tagB] at ($(ua)+(-4mm,0)$) {1};
  \node[anchor=north west,text width=79mm,align=left] at (100,-2) {
    \tracehead{P2Blue}{FROZEN STRUCTURE}\tracetitle{框架先定住，具体内容跟着真实任务变化}
    \tracebody{Brainstorm / Decision、Interaction Evidence、Technical Evidence、Final Confirmation、Global Evolution和Human--AI reflection各自有位置，但具体页数不预先限定。}
    \vspace{7mm}
    \tracehead{P2VioletAccent}{TRACEABILITY}\tracetitle{正文负责讲清过程，Evidence负责让我能回到原始记录}
    \tracebody{老师在正式页里先看到最关键的判断；需要追溯时，还能回到原始截图、Technical Evidence和最终确认。}
    \vspace{7mm}
    \tracehead{P2RoseAccent}{HANDOFF TO TASK 1}\tracetitle{下一阶段不再“设计证据”，而是在Task 1里自然产生证据}
    \tracebody{从这里开始，Experiment Decision Process只能来自Task 1里真实发生的方法讨论、参数选择、实验和我的判断，再交回Evidence Master整理。}
  };
\end{tikzpicture}
\provenance{本页使用真实“Process Report结构已成熟”的总结和我的最终确认。它标志Workflow Construction的预实验阶段完成，而不是Task 1研究结果。}
\end{document}
